% GENERATED FILE. Edit canonical lessons, then run npm run generate:books.
% canonical-source-hash: fnv1a64:c93bdd6dd05d549e
% canonical-lessons: PA-C273-tolna, PA-C273-takkri, PA-C273-killo, PA-C273-bhar, PA-C273-dhona-carry, PA-R273-review-267-273

\chapter{Two Kilos, Please}
\label{ch:pa-a2-two-kilos-please}
\hlchaptermodality{273}

\begin{chapteropening}
\textbf{By the end of this chapter:} \emph{I can ask for things by weight.}

Complete the last lesson of chapter 273: I can ask for things by weight.
\end{chapteropening}

\section[\texorpdfstring{\pa{ਤੋਲਣਾ} (tolṇā)}{ਤੋਲਣਾ (tolṇā)}]{\pa{ਤੋਲਣਾ} (tolṇā) — to weigh}

\label{lesson:PA-C273-tolna}



\begin{quote}
Before the new one: say the Punjabi for a rickshaw, then the Punjabi for a park.
\end{quote}

\subsection*{You'll want to know: \pa{ਤੋਲਣਾ}}
\textbf{\pa{ਤੋਲਣਾ}} — \emph{tolṇā} — \textquotedblleft{}to weigh\textquotedblright{}.

\begin{quote}
\pa{ਦੁਕਾਨ} \pa{ਵਿੱਚ} \pa{ਮੁੰਡਾ} \pa{ਆਲੂ} \pa{ਤੋਲਦਾ} \pa{ਹੈ।}
\end{quote}

\begin{quote}
\emph{dukān vich muṇḍā ālū toldā hai.}
\end{quote}

\begin{quote}
In the shop the boy weighs the potatoes.
\end{quote}

Past (I …): \pa{ਮੈਂ} \pa{ਆਲੂ} \pa{ਤੋਲੇ}, \emph{maiṁ ālū tole}.

\begin{cousinweb}
\textbf{\pa{ਤੋਲਣਾ}} (\emph{tolṇā}) is inherited from Sanskrit \emph{tolayati}, \textquotedblleft{}lifts, weighs\textquotedblright{}; CDIAL lists Punjabi \emph{toḷnā}, \textquotedblleft{}to weigh\textquotedblright{}, under it (CDIAL 5979).
\end{cousinweb}

\begin{grammarlens}[title={What you've built}]
One more word for this chapter.
\end{grammarlens}

\subsection*{Your turn}
\begin{itemize}
\raggedright
  \item \emph{Say it:} \emph{tolṇā}
  \item \emph{Say it:} \emph{tolṇā}, once more
  \item \emph{Say it:} \emph{dukān vich muṇḍā ālū toldā hai.} — In the shop the boy weighs the potatoes.
  \item \emph{Recall:} say the Punjabi for a rickshaw, then the Punjabi for a park, then say \emph{tolṇā} again
\end{itemize}

\begin{tcolorbox}[breakable,colback=teal!4,colframe=teal!35!black,title={Before you move on}]
What does \pa{ਤੋਲਣਾ} mean? (\textquotedblleft{}To weigh\textquotedblright{}.) Say it once more.
\end{tcolorbox}
\section[\texorpdfstring{\pa{ਤੱਕੜੀ} (takkṛī)}{ਤੱਕੜੀ (takkṛī)}]{\pa{ਤੱਕੜੀ} (takkṛī) — a pair of scales}

\label{lesson:PA-C273-takkri}



\begin{quote}
Before the new one: say the Punjabi for a park, then the Punjabi for to weigh.
\end{quote}

\subsection*{You'll want to know: \pa{ਤੱਕੜੀ}}
\textbf{\pa{ਤੱਕੜੀ}} — \emph{takkṛī} — \textquotedblleft{}a pair of scales\textquotedblright{}.

\begin{quote}
\pa{ਦੁਕਾਨ} \pa{ਦੀ} \pa{ਤੱਕੜੀ} \pa{ਪੁਰਾਣੀ} \pa{ਹੈ।}
\end{quote}

\begin{quote}
\emph{dukān dī takkṛī purāṇī hai.}
\end{quote}

\begin{quote}
The shop's scales are old.
\end{quote}

A feminine noun.

\begin{cousinweb}
\textbf{\pa{ਤੱਕੜੀ}} (\emph{takkṛī}) is filed by CDIAL under Sanskrit \emph{tarka}, \textquotedblleft{}inquiry, conjecture\textquotedblright{}; CDIAL lists Punjabi \emph{takkaṛī}, \textquotedblleft{}pair of scales\textquotedblright{}, there (CDIAL 5714).
\end{cousinweb}

\begin{grammarlens}[title={What you've built}]
One more word for this chapter.
\end{grammarlens}

\subsection*{Your turn}
\begin{itemize}
\raggedright
  \item \emph{Say it:} \emph{takkṛī}
  \item \emph{Say it:} \emph{takkṛī}, once more
  \item \emph{Say it:} \emph{dukān dī takkṛī purāṇī hai.} — The shop's scales are old.
  \item \emph{Recall:} say the Punjabi for a park, then the Punjabi for to weigh, then say \emph{takkṛī} again
  \item \emph{Write it:} copy \textbf{\pa{ਤੱਕੜੀ}} once, checking the whole word
\end{itemize}

\begin{tcolorbox}[breakable,colback=teal!4,colframe=teal!35!black,title={Before you move on}]
What does \pa{ਤੱਕੜੀ} mean? (\textquotedblleft{}A pair of scales\textquotedblright{}.) Say it once more.
\end{tcolorbox}
\section[\texorpdfstring{\pa{ਕਿੱਲੋ} (killo)}{ਕਿੱਲੋ (killo)}]{\pa{ਕਿੱਲੋ} (killo) — a kilo}

\label{lesson:PA-C273-killo}



\begin{quote}
Before the new one: say the Punjabi for to weigh, then the Punjabi for a pair of scales.
\end{quote}

\subsection*{You'll want to know: \pa{ਕਿੱਲੋ}}
\textbf{\pa{ਕਿੱਲੋ}} — \emph{killo} — \textquotedblleft{}a kilo\textquotedblright{}.

\begin{quote}
\pa{ਮੈਂ} \pa{ਦੋ} \pa{ਕਿੱਲੋ} \pa{ਸੇਬ} \pa{ਖ਼ਰੀਦਾਂਗਾ।}
\end{quote}

\begin{quote}
\emph{maiṁ do killo seb kharīdāṅgā.}
\end{quote}

\begin{quote}
I will buy two kilos of apples.
\end{quote}

A masculine noun.

\begin{cousinweb}
\textbf{\pa{ਕਿੱਲੋ}} (\emph{killo}) is borrowed from English \emph{kilo}.
\end{cousinweb}

\begin{grammarlens}[title={What you've built}]
One more word for this chapter.
\end{grammarlens}

\subsection*{Your turn}
\begin{itemize}
\raggedright
  \item \emph{Say it:} \emph{killo}
  \item \emph{Say it:} \emph{killo}, once more
  \item \emph{Say it:} \emph{maiṁ do killo seb kharīdāṅgā.} — I will buy two kilos of apples.
  \item \emph{Recall:} say the Punjabi for to weigh, then the Punjabi for a pair of scales, then say \emph{killo} again
  \item \emph{Write it:} copy \textbf{\pa{ਕਿੱਲੋ}} once, checking the whole word
\end{itemize}

\begin{tcolorbox}[breakable,colback=teal!4,colframe=teal!35!black,title={Before you move on}]
What does \pa{ਕਿੱਲੋ} mean? (\textquotedblleft{}A kilo\textquotedblright{}.) Say it once more.
\end{tcolorbox}
\section[\texorpdfstring{\pa{ਭਾਰ} (bhār)}{ਭਾਰ (bhār)}]{\pa{ਭਾਰ} (bhār) — weight; a load}

\label{lesson:PA-C273-bhar}



\begin{quote}
Before the new one: say the Punjabi for a pair of scales, then the Punjabi for a kilo.
\end{quote}

\subsection*{You'll want to know: \pa{ਭਾਰ}}
\textbf{\pa{ਭਾਰ}} — \emph{bhār} — \textquotedblleft{}weight; a load\textquotedblright{}.

\begin{quote}
\pa{ਆਲੂਆਂ} \pa{ਦਾ} \pa{ਭਾਰ} \pa{ਪੰਜ} \pa{ਕਿੱਲੋ} \pa{ਹੈ।}
\end{quote}

\begin{quote}
\emph{ālūāṁ dā bhār panj killo hai.}
\end{quote}

\begin{quote}
The potatoes weigh five kilos.
\end{quote}

A masculine noun.

\begin{cousinweb}
\textbf{\pa{ਭਾਰ}} (\emph{bhār}) is inherited from Sanskrit \emph{bhāra}, \textquotedblleft{}burden, weight\textquotedblright{}; CDIAL lists Punjabi \emph{bhār}, \textquotedblleft{}load\textquotedblright{}, under it (CDIAL 9459).
\end{cousinweb}

\begin{grammarlens}[title={What you've built}]
One more word for this chapter.
\end{grammarlens}

\subsection*{Your turn}
\begin{itemize}
\raggedright
  \item \emph{Say it:} \emph{bhār}
  \item \emph{Say it:} \emph{bhār}, once more
  \item \emph{Say it:} \emph{ālūāṁ dā bhār panj killo hai.} — The potatoes weigh five kilos.
  \item \emph{Recall:} say the Punjabi for a pair of scales, then the Punjabi for a kilo, then say \emph{bhār} again
  \item \emph{Write it:} copy \textbf{\pa{ਭਾਰ}} once, checking the whole word
\end{itemize}

\begin{tcolorbox}[breakable,colback=teal!4,colframe=teal!35!black,title={Before you move on}]
What does \pa{ਭਾਰ} mean? (\textquotedblleft{}Weight; a load\textquotedblright{}.) Say it once more.
\end{tcolorbox}
\section[\texorpdfstring{\pa{ਢੋਣਾ} (ḍhoṇā)}{ਢੋਣਾ (ḍhoṇā)}]{\pa{ਢੋਣਾ} (ḍhoṇā) — to carry a load}

\label{lesson:PA-C273-dhona-carry}



\begin{quote}
Before the new one: say the Punjabi for a kilo, then the Punjabi for weight.
\end{quote}

\subsection*{You'll want to know: \pa{ਢੋਣਾ}}
\textbf{\pa{ਢੋਣਾ}} — \emph{ḍhoṇā} — \textquotedblleft{}to carry a load\textquotedblright{}.

\begin{quote}
\pa{ਮੈਂ} \pa{ਬਾਜ਼ਾਰ} \pa{ਤੋਂ} \pa{ਸਬਜ਼ੀ} \pa{ਘਰ} \pa{ਢੋਂਦਾ} \pa{ਹਾਂ।}
\end{quote}

\begin{quote}
\emph{maiṁ bāzār ton sabzī ghar ḍhondā hāṁ.}
\end{quote}

\begin{quote}
I carry the vegetables home from the market.
\end{quote}

Past (I …): \pa{ਮੈਂ} \pa{ਆਲੂ} \pa{ਢੋਏ}, \emph{maiṁ ālū ḍhoe}.

\begin{cousinweb}
\textbf{\pa{ਢੋਣਾ}} (\emph{ḍhoṇā}) is inherited from Sanskrit \emph{ḍhaukayati}, \textquotedblleft{}brings near\textquotedblright{}; CDIAL lists Punjabi \emph{ḍhoṇā}, \textquotedblleft{}to carry, convey\textquotedblright{}, under it (CDIAL 5610).
\end{cousinweb}

\begin{grammarlens}[title={What you've built}]
One more word for this chapter.
\end{grammarlens}

\subsection*{Your turn}
\begin{itemize}
\raggedright
  \item \emph{Say it:} \emph{ḍhoṇā}
  \item \emph{Say it:} \emph{ḍhoṇā}, once more
  \item \emph{Say it:} \emph{maiṁ bāzār ton sabzī ghar ḍhondā hāṁ.} — I carry the vegetables home from the market.
  \item \emph{Recall:} say the Punjabi for a kilo, then the Punjabi for weight, then say \emph{ḍhoṇā} again
  \item \emph{Write it:} copy \textbf{\pa{ਢੋਣਾ}} once, checking the whole word
\end{itemize}

\begin{tcolorbox}[breakable,colback=teal!4,colframe=teal!35!black,title={Before you move on}]
What does \pa{ਢੋਣਾ} mean? (\textquotedblleft{}To carry a load\textquotedblright{}.) Say it once more.
\end{tcolorbox}
\section[(dialogue)]{Looking back: from Which Floor, Which Office? to Two Kilos, Please}

\label{lesson:PA-R273-review-267-273}



\begin{quote}
Say the Punjabi for a police station and a storey.
\end{quote}

\subsection*{Your turn}
Say these aloud:

\begin{itemize}
\raggedright
  \item a police station, a storey, a post office, a court of law, whereabouts
  \item a bus stand, a bus conductor, a seat, a crowd, to stay put
  \item to recognize, a map, a mark, to draw near, a pothole
  \item to search for, to stand, to see someone across, a mile, open ground
  \item north, south, east, west, to show someone around
  \item to get separated, a police constable, a police post, a rickshaw, a park
  \item to weigh, a pair of scales, a kilo, weight, to carry a load
\end{itemize}

\begin{tcolorbox}[breakable,colback=teal!4,colframe=teal!35!black,title={Before you move on}]
A police station? (\textbf{\pa{ਥਾਣਾ}}, \emph{thāṇā}.) A storey? (\textbf{\pa{ਮੰਜ਼ਿਲ}}, \emph{manzil}.) A post office? (\textbf{\pa{ਡਾਕਖ਼ਾਨਾ}}, \emph{ḍākkhānā}.) A court of law? (\textbf{\pa{ਅਦਾਲਤ}}, \emph{adālat}.) Whereabouts? (\textbf{\pa{ਠਿਕਾਣਾ}}, \emph{ṭhikāṇā}.) A bus stand? (\textbf{\pa{ਅੱਡਾ}}, \emph{aḍḍā}.) A bus conductor? (\textbf{\pa{ਕੰਡਕਟਰ}}, \emph{kaṇḍakṭar}.) A seat? (\textbf{\pa{ਸੀਟ}}, \emph{sīṭ}.) A crowd? (\textbf{\pa{ਭੀੜ}}, \emph{bhīṛ}.) To stay put? (\textbf{\pa{ਟਿਕਣਾ}}, \emph{ṭikṇā}.) To recognize? (\textbf{\pa{ਪਛਾਣਨਾ}}, \emph{pachhāṇnā}.) A map? (\textbf{\pa{ਨਕਸ਼ਾ}}, \emph{nakshā}.) A mark? (\textbf{\pa{ਨਿਸ਼ਾਨ}}, \emph{nishān}.) To draw near? (\textbf{\pa{ਢੁੱਕਣਾ}}, \emph{ḍhukkṇā}.) A pothole? (\textbf{\pa{ਟੋਆ}}, \emph{ṭoā}.) To search for? (\textbf{\pa{ਢੂੰਢਣਾ}}, \emph{ḍhūṇḍhṇā}.) To stand? (\textbf{\pa{ਖੜ੍ਹਨਾ}}, \emph{khaṛhnā}.) To see someone across? (\textbf{\pa{ਲੰਘਾਉਣਾ}}, \emph{laṅghāuṇā}.) A mile? (\textbf{\pa{ਮੀਲ}}, \emph{mīl}.) Open ground? (\textbf{\pa{ਮੈਦਾਨ}}, \emph{maidān}.) North? (\textbf{\pa{ਉੱਤਰ}}, \emph{uttar}.) South? (\textbf{\pa{ਦੱਖਣ}}, \emph{dakkhaṇ}.) East? (\textbf{\pa{ਪੂਰਬ}}, \emph{pūrab}.) West? (\textbf{\pa{ਪੱਛਮ}}, \emph{pachchham}.) To show someone around? (\textbf{\pa{ਘੁਮਾਉਣਾ}}, \emph{ghumāuṇā}.) To get separated? (\textbf{\pa{ਵਿਛੜਨਾ}}, \emph{vichhaṛnā}.) A police constable? (\textbf{\pa{ਸਿਪਾਹੀ}}, \emph{sipāhī}.) A police post? (\textbf{\pa{ਚੌਕੀ}}, \emph{chaukī}.) A rickshaw? (\textbf{\pa{ਰਿਕਸ਼ਾ}}, \emph{rikshā}.) A park? (\textbf{\pa{ਪਾਰਕ}}, \emph{pārak}.) To weigh? (\textbf{\pa{ਤੋਲਣਾ}}, \emph{tolṇā}.) A pair of scales? (\textbf{\pa{ਤੱਕੜੀ}}, \emph{takkṛī}.) A kilo? (\textbf{\pa{ਕਿੱਲੋ}}, \emph{killo}.) Weight? (\textbf{\pa{ਭਾਰ}}, \emph{bhār}.) To carry a load? (\textbf{\pa{ਢੋਣਾ}}, \emph{ḍhoṇā}.)
\end{tcolorbox}
